\section{Spectral sequences: double complexes}
\label{section-double-complex}

\noindent
Let $K^{\bullet, \bullet}$ be a double complex, see
Section \ref{section-double-complexes}. It is customary to denote
$H^p_I(K^{\bullet, \bullet})$
the complex with terms $\Ker(d_1^{p, q})/\Im(d_1^{p - 1, q})$
(varying $q$) and differential induced by $d_2$.
Then $H^q_{II}(H^p_I(K^{\bullet, \bullet}))$ denotes its cohomology in
degree $q$. It is also customary to denote $H^q_{II}(K^{\bullet, \bullet})$
the complex with terms $\Ker(d_2^{p, q})/\Im(d_2^{p, q - 1})$
(varying $p$) and differential induced by $d_1$.
Then $H^p_I(H^q_{II}(K^{\bullet, \bullet}))$ denotes its cohomology in
degree $p$. It will turn out that these cohomology groups show up
as the terms in the spectral sequence for a filtration on the
associated total complex or simple complex, see
Definition \ref{definition-associated-simple-complex}.

\medskip\noindent
There are two natural filtrations on the total complex
$\text{Tot}(K^{\bullet, \bullet})$
associated to the double complex $K^{\bullet, \bullet}$. Namely, we
define
$$
F_I^p(\text{Tot}^n(K^{\bullet, \bullet})) =
\bigoplus\nolimits_{i + j = n, \ i \geq p} K^{i, j}
\quad
\text{and}
\quad
F_{II}^p(\text{Tot}^n(K^{\bullet, \bullet})) =
\bigoplus\nolimits_{i + j = n, \ j \geq p} K^{i, j}.
$$
It is immediately verified that $(\text{Tot}(K^{\bullet, \bullet}), F_I)$ and
$(\text{Tot}(K^{\bullet, \bullet}), F_{II})$ are filtered complexes.
By Section \ref{section-filtered-complex}
we obtain two spectral sequences. It is customary to
denote $({}'E_r, {}'d_r)_{r \geq 0}$ the spectral sequence associated
to the filtration $F_I$ and to denote $({}''E_r, {}''d_r)_{r \geq 0}$
the spectral sequence associated to the filtration $F_{II}$.
Here is a description of these spectral sequences.

\begin{lemma}
\label{lemma-ss-double-complex}
Let $\mathcal{A}$ be an abelian category.
Let $K^{\bullet, \bullet}$ be a double complex.
The spectral sequences associated to $K^{\bullet, \bullet}$
have the following terms:
\begin{enumerate}
\item ${}'E_0^{p, q} = K^{p, q}$ with
${}'d_0^{p, q} = (-1)^p d_2^{p, q} : K^{p, q} \to K^{p, q + 1}$,
\item ${}''E_0^{p, q} = K^{q, p}$ with
${}''d_0^{p, q} = d_1^{q, p} : K^{q, p} \to K^{q + 1, p}$,
\item ${}'E_1^{p, q} = H^q(K^{p, \bullet})$ with
${}'d_1^{p, q} = H^q(d_1^{p, \bullet})$,
\item ${}''E_1^{p, q} = H^q(K^{\bullet, p})$ with
${}''d_1^{p, q} = (-1)^q H^q(d_2^{\bullet, p})$,
\item ${}'E_2^{p, q} = H^p_I(H^q_{II}(K^{\bullet, \bullet}))$,
\item ${}''E_2^{p, q} = H^p_{II}(H^q_I(K^{\bullet, \bullet}))$.
\end{enumerate}
\end{lemma}

\begin{proof}
Omitted.
\end{proof}

\noindent
These spectral sequences define two filtrations on
$H^n(\text{Tot}(K^{\bullet, \bullet}))$.
We will denote these $F_I$ and $F_{II}$.

\begin{definition}
\label{definition-ss-double-complex-converge}
Let $\mathcal{A}$ be an abelian category.
Let $K^{\bullet, \bullet}$ be a double complex.
We say the spectral sequence $({}'E_r, {}'d_r)_{r \geq 0}$
{\it weakly converges to $H^n(\text{Tot}(K^{\bullet, \bullet}))$},
{\it abuts to $H^n(\text{Tot}(K^{\bullet, \bullet}))$}, or
{\it converges to $H^n(\text{Tot}(K^{\bullet, \bullet}))$}
if Definition \ref{definition-filtered-complex-ss-converges} applies.
Similarly we say the spectral sequence $({}''E_r, {}''d_r)_{r \geq 0}$
{\it weakly converges to $H^n(\text{Tot}(K^{\bullet, \bullet}))$},
{\it abuts to $H^n(\text{Tot}(K^{\bullet, \bullet}))$}, or
{\it converges to $H^n(\text{Tot}(K^{\bullet, \bullet}))$}
if Definition \ref{definition-filtered-complex-ss-converges} applies.
\end{definition}

\noindent
As mentioned above there is no consistent terminology regarding these
notions in the literature. In the situation of the definition, we have
weak convergence of the first spectral sequence if for all $n$
$$
\text{gr}_{F_I}(H^n(\text{Tot}(K^{\bullet, \bullet}))) =
\oplus_{p + q = n} {}'E_\infty^{p, q}
$$
via the canonical comparison of
Lemma \ref{lemma-compute-cohomology-filtered-complex}.
Similarly the second spectral sequence $({}''E_r, {}''d_r)_{r \geq 0}$
weakly converges if for all $n$
$$
\text{gr}_{F_{II}}(H^n(\text{Tot}(K^{\bullet, \bullet}))) =
\oplus_{p + q = n} {}''E_\infty^{p, q}
$$
via the canonical comparison of
Lemma \ref{lemma-compute-cohomology-filtered-complex}.

\begin{lemma}
\label{lemma-first-quadrant-ss}
Let $\mathcal{A}$ be an abelian category. Let $K^{\bullet, \bullet}$
be a double complex. Assume that for every $n \in \mathbf{Z}$ there are
only finitely many nonzero $K^{p, q}$ with $p + q = n$. Then
\begin{enumerate}
\item the two spectral sequences associated to $K^{\bullet, \bullet}$
are bounded,
\item the filtrations $F_I$, $F_{II}$ on each
$H^n(\text{Tot}(K^{\bullet, \bullet}))$ are finite,
\item the spectral sequences $({}'E_r, {}'d_r)_{r \geq 0}$ and
$({}''E_r, {}''d_r)_{r \geq 0}$ converge to
$H^*(\text{Tot}(K^{\bullet, \bullet}))$,
\item if $\mathcal{C} \subset \mathcal{A}$ is a weak Serre subcategory
and for some $r$ we have ${}'E_r^{p, q} \in \mathcal{C}$ for all
$p, q \in \mathbf{Z}$, then $H^n(\text{Tot}(K^{\bullet, \bullet}))$
is in $\mathcal{C}$. Similarly for $({}''E_r, {}''d_r)_{r \geq 0}$.
\end{enumerate}
\end{lemma}

\begin{proof}
Follows immediately from Lemma \ref{lemma-biregular-ss-converges}.
\end{proof}

\noindent
Here is our first application of spectral sequences.

\begin{lemma}
\label{lemma-double-complex-gives-resolution}
Let $\mathcal{A}$ be an abelian category.
Let $K^\bullet$ be a complex.
Let $A^{\bullet, \bullet}$ be a double complex.
Let $\alpha^p : K^p \to A^{p, 0}$ be morphisms.
Assume that
\begin{enumerate}
\item For every $n \in \mathbf{Z}$ there are only finitely many nonzero
$A^{p, q}$ with $p + q = n$.
\item We have $A^{p, q} = 0$ if $q < 0$.
\item The morphisms $\alpha^p$ give rise to a morphism
of complexes $\alpha : K^\bullet \to A^{\bullet, 0}$.
\item The complex $A^{p, \bullet}$ is exact in all degrees
$q \not = 0$ and the morphism $K^p \to A^{p, 0}$ induces
an isomorphism $K^p \to \Ker(d_2^{p, 0})$.
\end{enumerate}
Then $\alpha$ induces a quasi-isomorphism
$$
K^\bullet \longrightarrow \text{Tot}(A^{\bullet, \bullet})
$$
of complexes.
Moreover, there is a variant of this lemma involving the second
variable $q$ instead of $p$.
\end{lemma}

\begin{proof}
The map is simply the map given by the morphisms
$K^n \to A^{n, 0} \to \text{Tot}^n(A^{\bullet, \bullet})$,
which are easily seen to define a morphism of complexes.
Consider the spectral sequence $({}'E_r, {}'d_r)_{r \geq 0}$
associated to the double complex $A^{\bullet, \bullet}$.
By Lemma \ref{lemma-first-quadrant-ss} this spectral sequence converges
and the induced filtration on $H^n(\text{Tot}(A^{\bullet, \bullet}))$
is finite for each $n$.
By Lemma \ref{lemma-ss-double-complex} and assumption (4) we have
${}'E_1^{p, q} = 0$ unless $q = 0$ and $'E_1^{p, 0} = K^p$
with differential ${}'d_1^{p, 0}$ identified with $d_K^p$.
Hence ${}'E_2^{p, 0} = H^p(K^\bullet)$ and zero otherwise.
This clearly implies ${}'d_2^{p, q} = {}'d_3^{p, q} = \ldots = 0$
for degree reasons. Hence we conclude that
$H^n(\text{Tot}(A^{\bullet, \bullet})) = H^n(K^\bullet)$.
We omit the verification that this identification is given by the
morphism of complexes $K^\bullet \to \text{Tot}(A^{\bullet, \bullet})$
introduced above.
\end{proof}

\begin{lemma}
\label{lemma-homotopy-complex-complexes}
Let $\mathcal{A}$ be an abelian category.
Let $M^\bullet$ be a complex of $\mathcal{A}$. Let
$$
a :
M^\bullet[0]
\longrightarrow
\left(A^{0, \bullet} \to A^{1, \bullet} \to A^{2, \bullet} \to \ldots \right)
$$
be a homotopy equivalence in the category of complexes of complexes
of $\mathcal{A}$. Then the map
$\alpha : M^\bullet \to \text{Tot}(A^{\bullet, \bullet})$
induced by $M^\bullet \to A^{0, \bullet}$ is a homotopy equivalence.
\end{lemma}

\begin{proof}
The statement makes sense as a complex of complexes is the same thing
as a double complex. The assumption means there is a map
$$
b :
\left(A^{0, \bullet} \to A^{1, \bullet} \to A^{2, \bullet} \to \ldots \right)
\longrightarrow
M^\bullet[0]
$$
such that $a \circ b$ and $b \circ a$ are homotopic to the identity
in the category of complexes of complexes. This means that $b \circ a$
is the identity of $M^\bullet[0]$ (because there is only one term in
degree $0$). Also, observe that $b$ is given by a map
$b^0 : A^{0, \bullet} \to M^\bullet$ and zero in all other degrees.
Thus $b$ induces a map
$\beta : \text{Tot}(A^{\bullet, \bullet}) \to M^\bullet$
and $\beta \circ \alpha$ is the identity on $M^\bullet$.
Finally, we have to show that the map
$\alpha \circ \beta$ is homotopic to the identity.
For this we choose maps of complexes
$h^n : A^{n, \bullet} \to A^{n - 1, \bullet}$
such that $a \circ b - \text{id} = d_1 \circ h + h \circ d_1$
which exist by assumption. Here $d_1 : A^{n, \bullet} \to A^{n + 1, \bullet}$
are the differentials of the complex of complexes. We will also
denote $d_2$ the differentials of the complexes $A^{n, \bullet}$
for all $n$. Let $h^{n, m} : A^{n, m} \to A^{n - 1, m}$ be the components of
$h^n$. Then we can consider
$$
h' : \text{Tot}(A^{\bullet, \bullet})^k =
\bigoplus\nolimits_{n + m = k} A^{n, m}
\to
\bigoplus\nolimits_{n + m = k - 1} A^{n, m} =
\text{Tot}(A^{\bullet, \bullet})^{k - 1}
$$
given by $h^{n, m}$ on the summand $A^{n, m}$. Then we compute
that the map
$$
d_{\text{Tot}(A^{\bullet, \bullet})} \circ h' +
h' \circ d_{\text{Tot}(A^{\bullet, \bullet})}
$$
restricted to the summand $A^{n, m}$ is equal to
$$
d_1^{n - 1, m} \circ h^{n, m} +
(-1)^{n - 1} d_2^{n - 1, m} \circ h^{n, m} +
h^{n + 1, m} \circ d_1^{n, m} + h^{n, m + 1} \circ (-1)^nd_2^{n, m}
$$
Since $h^n$ is a map of complexes, the terms
$(-1)^{n - 1} d_2^{n - 1, m} \circ h^{n, m}$ and
$h^{n, m + 1} \circ (-1)^nd_2^{n, m}$ cancel.
The other two terms give
$(\alpha \circ \beta)|_{A^{n, m}} - \text{id}_{A^{n, m}}$
because $a \circ b - \text{id} = d_1 \circ h + h \circ d_1$.
This finishes the proof.
\end{proof}








